\chapter{Implementing threads}
\label{ch:impl}

Chapter~2 used threads; this chapter builds them. The mechanics are the
same whether the switching happens in a user-level library
(\module{10}) or in a kernel (SWEB): a record per thread, a way to save
and restore registers, a stack per thread, a state machine, and wait
queues. The hard parts are at the edges --- how a thread starts, how it
ends, who frees its stack, and what happens when an interrupt arrives in
the middle of the library's own bookkeeping.

\section{The thread control block and the context}

For each thread the implementation keeps a \emph{thread control block}
(TCB): id, state, saved registers (the \emph{context}), stack, start
routine and argument, return value, detached flag, and links for the
queues it may sit in.

A context switch saves the running thread's registers into its TCB and
loads the next thread's registers from its TCB. Loading the stack pointer
switches stacks; loading the instruction pointer continues wherever the
next thread stopped. From the thread's point of view, the call that
switched away simply returns --- later. \code{swapcontext(\&old, \&new)}
does exactly this in user space; SWEB's \code{arch\_contextSwitch}
restores a \code{Thread}'s saved \code{ArchThreadRegisters}.

\section{States}

\begin{figure}[htbp]
\centering
\begin{tikzpicture}[font=\small,>=Latex,node distance=3.4cm,
    st/.style={draw,rounded corners=3pt,minimum width=1.8cm,minimum height=0.75cm,fill=newbg}]
  \node[st] (ready) {READY};
  \node[st,right=of ready] (run) {RUNNING};
  \node[st,below=2.2cm of $(ready)!0.5!(run)$] (blk) {BLOCKED};
  \node[st,right=2.2cm of run,fill=warnbg] (zom) {ZOMBIE};
  \node[st,right=1.8cm of zom,fill=black!8] (free) {freed};
  \node[left=1.2cm of ready] (new) {};
  \draw[->] (new) -- node[above] {create} (ready);
  \draw[->] (ready) to[bend left=25] node[above] {scheduled} (run);
  \draw[->] (run) to[bend left=25] node[below,fill=white,inner sep=1pt] {yield / preempted} (ready);
  \draw[->] (run.south) -- node[right,pos=0.55] {join, lock, wait} (blk.north east);
  \draw[->] (blk.north west) -- node[left,pos=0.45] {woken} (ready.south);
  \draw[->] (run) -- node[above] {exit} (zom);
  \draw[->] (zom) -- node[above,align=center] {joined /\\detached} (free);
\end{tikzpicture}
\caption{Thread states. SWEB's names: \code{Running} covers READY and
RUNNING, \code{Sleeping} is BLOCKED, \code{ToBeDestroyed} is the way out.}
\label{fig:states}
\end{figure}

A BLOCKED thread is in no ready queue; exactly one wait queue (a mutex's,
a condition's, a joiner slot) knows it, and whoever makes the awaited
event happen moves it back to READY.

\section{Birth: a stack and a trampoline}

A new thread needs a stack of its own and a first context that ``returns''
into its code. The first instruction must not be the start routine itself:
when \code{start\_routine} returns, it would return to whatever garbage
lies on the fresh stack. Instead the thread starts in a small
\emph{trampoline}:

\begin{lstlisting}
static void trampoline(void) {
  after_switch();                                 /* see "Death" below */
  uthread_exit(current->start(current->arg));      /* return == exit */
}
\end{lstlisting}

Stacks should have a \emph{guard page} below them (mapped without any
access rights), so that an overflow faults at once instead of silently
overwriting the neighbouring stack.

\section{Death: exit, zombies, join, detach}

\begin{itemize}
  \item \textbf{exit}: store the return value, become a ZOMBIE, wake the
    joiner if one waits, switch away --- forever.
  \item \textbf{join}: if the target is already a zombie, collect and free
    it; otherwise block until its exit wakes you, then collect and free.
    Errors: no such thread (\code{ESRCH}), joining yourself or closing a
    cycle of joins (\code{EDEADLK}), a detached target or a second joiner
    (\code{EINVAL}).
  \item \textbf{detach}: nobody will join; free the thread as soon as it
    is a zombie (at once, if it already is).
\end{itemize}

\paragraph{Nobody frees their own stack.} A detached thread that exits
cannot unmap its stack --- it is still running on it; the next push or
return would fault. The freeing has to happen \emph{after} the switch, by
somebody else: \module{10}'s library notes the thread in a variable, and
the next thread to run frees it right after the switch (in
\code{after\_switch}, which the trampoline also calls). SWEB solves the
same problem for kernel stacks with a dedicated thread:
\code{Thread::kill} only marks the thread \code{ToBeDestroyed}, and the
cleanup thread \code{delete}s it later.

\section{Blocking primitives}

A blocking mutex is an owner plus a FIFO of waiting TCBs:
\begin{itemize}
  \item lock: free $\to$ take it; else enqueue yourself, state BLOCKED,
    switch away;
  \item unlock: waiters? either \emph{hand the mutex over} to the first
    one (FIFO, fair, but every contended unlock forces a switch) or free
    it and wake a waiter who must compete again (\emph{barging}: faster,
    unfair --- SWEB's \code{Mutex} does this).
\end{itemize}

A condition variable is just a wait queue; \code{wait} enqueues,
unlocks the mutex and blocks \emph{without any switch in between},
then re-locks after being woken. Deadlock detection comes almost for free:
if the ready queue is empty but some thread is BLOCKED, nothing can ever
wake it.

\section{Preemption and interrupt safety}

Without preemption, a thread runs until it calls into the library;
atomicity inside the library is automatic. With a timer
(\module{10}'s bonus uses \code{SIGVTALRM}, the kernel uses the timer
interrupt), a switch can happen \emph{anywhere} --- also in the middle of
\code{q\_push} or between ``mutex is free'' and ``I am the owner''. The
library's own data then needs protection: block the signal (user space)
or disable interrupts (kernel, one CPU) around every operation on it,
exactly as SWEB's scheduler runs with interrupts off. On several CPUs this
is not enough; a spinlock is needed as well.

And code that can be interrupted at any instruction must not be
re-entered through non-reentrant functions: a thread preempted inside
\code{malloc} holds malloc's lock; another thread on the same kernel
thread calling \code{malloc} deadlocks. SWEB's version of the rule: no
\code{new}/\code{delete} with interrupts off or in interrupt handlers.

\section{User-level vs.\ kernel-level, revisited}

\begin{center}
\begin{tabular}{L{3.6cm}L{4.8cm}L{4.8cm}}
\toprule
 & \textbf{User-level (M:1)} & \textbf{Kernel-level (1:1)} \\
\midrule
switch cost & function call + registers (\module{10}: $\approx$0.4 µs,
mostly glibc's signal-mask syscall) & kernel entry, scheduler, exit
($\approx$3 µs in \module{10}'s measurement) \\
blocking syscall & blocks \emph{all} threads & blocks only the caller \\
multiple cores & no & yes \\
per thread & TCB, user stack & kernel stack, user stack, two register
sets, kernel TCB \\
\bottomrule
\end{tabular}
\end{center}

\begin{swebbox}
After A1 every pthread is a kernel \code{Thread} (1:1). The pieces map
one to one:
\begin{itemize}
  \item TCB = \code{Thread}; its \code{kernel\_stack\_} is 8~KiB inside
    the object; \code{user\_registers\_} and \code{kernel\_registers\_} are
    the two contexts (user mode, and the kernel side of a system call or
    interrupt).
  \item Birth: \code{ArchThreads::createUserRegisters(regs, entry,
    user\_stack\_top, kernel\_stack)}; \code{entry} = a libc wrapper that
    calls \code{start\_routine(arg)} and then \code{pthread\_exit}; the two
    values reach it in \code{rdi}/\code{rsi}. \code{setAddressSpace} gives
    the thread the process's page tables.
  \item The user stack: a new region in the shared address space, below
    the first thread's stack page, with a guard gap (chapter~8).
  \item Death and reaping: \code{kill} $\to$ \code{ToBeDestroyed} $\to$
    cleanup thread. Join data (return value, finished) must survive the
    \code{Thread} object --- keep it in the process.
  \item Waiting: \code{Condition}/\code{Mutex} (chapter~5), never a yield
    loop (chapter~9).
\end{itemize}
\end{swebbox}

\begin{keybox}
\begin{itemize}
  \item A context switch = save registers into one TCB, load them from
    another; the stack pointer switches the stack.
  \item Start through a trampoline so that returning means exiting; guard
    pages catch stack overflows.
  \item Zombies keep the return value until join; detached threads are
    freed at exit --- by somebody else, after the switch.
  \item Blocking = a wait queue + state BLOCKED + switch; hand-off vs.
    barging.
  \item Preemption makes the library's own data shared: disable
    interrupts / block the signal; no non-reentrant calls.
\end{itemize}
\end{keybox}
